\section{Background}
\label{sec:background}

The Recursive Spawning Protocol (RSP) draws from five distinct research traditions: agent lifecycle management in multi-agent systems; accountability and provenance in distributed AI; fixed versus mutable delegation chains; hierarchical task decomposition; and the BX3 Framework as the minimal sufficient substrate for immutable parent chains. This section surveys each tradition, identifies the structural gaps RSP fills, and shows why those gaps cannot be closed by extensions of existing approaches.

% -------------------------------------------------------------------------%
\subsection{Agent Lifecycle Management in Multi-Agent Systems}
% -------------------------------------------------------------------------%

Governing agent populations across their full lifecycle---from provisioning through decommissioning---has become a central research concern as autonomous agents move from single-task deployments into complex, multi-agent workflows. The Cloud Native Computing Foundation's 2026 cloud-native agentic standards identify \emph{agent identity}, least-privilege access, and tenancy controls as foundational requirements for multi-agent systems in Kubernetes environments~\cite{cncf-agentic-standards}. Microsoft's Agent Governance Toolkit operationalizes these through Agent OS (a stateless policy engine intercepting agent actions at sub-millisecond latency) and Agent Mesh (cryptographic DIDs with Ed25519 signatures, an Inter-Agent Trust Protocol, and dynamic trust scoring)~\cite{microsoft-agent-governance}. Neither approach addresses the structural immutability of the parent delegation chain that RSP requires.

The Agent Lifecycle Protocol (ALP) defines a formal state machine governing six stages---Genesis, Fork, Migration, Retirement, Succession, and Decommission---with explicit transition rules and a genealogical registry~\cite{alp2025}. ALP identifies \emph{lineage tracking} and \emph{accountability across generations} as unsolved challenges, noting that a comprehensive registry for querying an agent's complete family and history is necessary for audit and compliance but absent from conventional runtimes. RSP directly addresses this: the immutable parent pointer $\alpha(n)$ is the foundational primitive for the genealogical registry ALP calls for, and the Fact Layer's forensic ledger provides the append-only chain of custody that makes lineage verification runtime-computable rather than reconstructed post-hoc from fragmented logs.

The Human Delegation Protocol (HDP) introduces JSON-based tokens binding a human principal to authorization scope and session, with each agent in a multi-hop chain appending signed hops to an append-only provenance record~\cite{hdp-draft,hdp-draft-2}. HDP demonstrates that multi-hop, append-only human provenance in agentic workflows is achievable and verifiable offline using only the issuer's public key. RSP adopts HDP's append-only chain principle and extends it structurally: while HDP achieves provenance through signed hop records, RSP achieves structural immutability of the parent pointer at the Fact Layer write protocol level, eliminating the possibility of retroactive chain modification regardless of the signing scheme employed.

The core structural problem RSP addresses is \emph{spawn chain opacity}: the inability to determine, at any point during or after execution, whether a given agent's authority traces to a human principal or to an autonomous predecessor that may have drifted from its authorized purpose. This opacity arises because parent pointers are mutable in conventional architectures---they can be redirected by administrative action, policy override, or runtime reconfiguration---making retroactive provenance reconstruction inherently unreliable.

% -------------------------------------------------------------------------%
\subsection{Accountability and Provenance in Distributed Systems}
% -------------------------------------------------------------------------%

The accountability problem in distributed systems---how to construct a reliable, tamper-evident record of which actor performed which action, under what authorization, and with what outcome---has been studied extensively, but AI agent chains introduce structural challenges that existing approaches do not fully address.

The W3C PROV data model provides a foundational framework for representing provenance information, defining \emph{entity}, \emph{activity}, and \emph{agent} with relationships (\texttt{wasGeneratedBy}, \texttt{wasAssociatedWith}, \texttt{wasAttributedTo}) that trace how data or artifacts were produced and by whom~\cite{moreau2015-prov}. PROV-AGENT extends PROV for AI agent-centric workflows, introducing agent artifacts (tools, prompts, responses, model invocations) as first-class PROV entities integrated with non-agentic tasks in a unified provenance graph~\cite{prov-agent2025}. However, PROV-AGENT's provenance graph remains a post-hoc record of observed events; it does not structurally enforce that every agent in the graph was provisioned through a scope-refined, depth-bounded chain from a human anchor.

Certificate Transparency demonstrates that append-only hash-chained logs can provide reliable audit trails even when participants behave adversarially~\cite{laurie2013certificate}. The Bitcoin blockchain extends this to decentralized, consensus-based append-only ledgers, maintaining tamper-evidence across untrusted network participants~\cite{nakamoto2008bitcoin}. RSP applies this principle to the agent spawn chain: the Fact Layer maintains a hash-chained forensic ledger in which every spawn event, state transition, and authorization directive is permanently recorded. Any attempt to insert, delete, or reorder a historical entry breaks hash continuity and is detectable by any observer with ledger read access.

SentinelAgent provides formal verification of immutable delegation chains in federal multi-agent AI systems, defining the Delegation Chain Calculus (DCC) with seven properties---authority narrowing, policy preservation, forensic reconstructibility, cascade containment, scope-action conformance, output schema conformance (six deterministic), and intent preservation (probabilistic)---alongside four meta-theorems including property minimality and graceful degradation under adversarial conditions~\cite{sentinelagent2026}. SentinelAgent's Intent-Preserving Delegation Protocol achieves high detection accuracy on DelegationBench v4, demonstrating that structural enforcement of delegation chain properties is achievable and verifiable.

The Agent Identity Protocol (AIP) introduces Invocation-Bound Capability Tokens (IBCTs) that fuse identity, attenuated authorization, and provenance into an append-only token chain for MCP and A2A communications, demonstrating sub-millisecond verification speeds and 100\% adversarial rejection rates~\cite{aip2025}.

The fundamental accountability challenge RSP addresses is that conventional provenance systems record what happened but do not structurally enforce that each step in the chain satisfied scope refinement and depth constraints at spawn time. RSP makes these constraints first-class, structurally enforced properties of the Fact Layer's write protocol, converting accountability from a post-hoc reconstructive discipline into a runtime verification property.

% -------------------------------------------------------------------------%
\subsection{Fixed and Mutable Delegation Chains}
% -------------------------------------------------------------------------%

The structural distinction between fixed and mutable delegation chains is the central design choice determining whether accountability is achievable in a multi-agent system. A mutable delegation chain allows the parent pointer $\alpha(n)$ to be rewritten after provisioning. A fixed (immutable) delegation chain establishes $\alpha(n)$ once at provisioning and enforces immutability through the Fact Layer write protocol, making it structurally impossible to redirect after the fact.

Mutable delegation chains are the historical norm. OAuth 2.1's token exchange model enables scoped, time-limited delegation but leaves the delegation graph itself mutable: authorized users can re-delegate, administrators can modify scope, and tokens can be revoked retroactively. UCAN (User Controlled Authorization Networks) introduces self-sovereign authorization with a capability model where delegation chains are encoded as capability graphs verified transitively; UCAN's capability graph is mutable, as capabilities can be delegated, transferred, and revoked with revocation propagating through the graph~\cite{UCAN}. This is appropriate for human-controlled authorization networks where the principal remains in the loop, but is structurally insufficient for autonomous agent chains where the human principal may not be reachable at decision time.

DToken provides a lightweight, traceable delegation architecture for distributed systems using token-based delegation with chain formation where each step records origin, delegator, and delegatee, enabling end-to-end provenance tracking~\cite{dtoken2025}. RFC-008's Delegated Authority Envelopes extend this with monotonically narrowing scope, parent-authority hash chaining, and PDP enforcement, forming cryptographically signed delegation chains where each link can only narrow its parent's authority~\cite{capisco-rfc008}.

The IETF draft on identity management for agentic AI identifies recursive delegation risk---agents creating sub-agents or delegating tasks introducing complex chains of accountability and control---as a critical future challenge, with preemptive authorization and careful scoping proposed as mitigations~\cite{openid-agentic-ai}. Google DeepMind's \emph{Intelligent AI Delegation} framework identifies recursive delegation risk as a critical failure mode in long agent chains, proposing that delegation should transfer scoped authority, responsibility, and accountability, with monitoring and trust mechanisms that hold when environments change or agents fail~\cite{intelligent-delegation}.

RSP provides the structural complement to these policy-level approaches: rather than relying on the correctness of individual authorization decisions, RSP enforces that every spawn event satisfies the scope subset constraint and depth bound as structural invariants. The fundamental problem with mutable delegation chains in multi-agent contexts is that they permit \emph{re-parenting drift}: a node provisioned under one human anchor can be retroactively re-parented to a different chain, severing the provenance link to the original principal and enabling actions the original principal never authorized. This failure mode is architecturally impossible in RSP because the Fact Layer write protocol rejects any modification to $\alpha(n)$ after provisioning.

% -------------------------------------------------------------------------%
\subsection{Hierarchical Task Decomposition in AI}
% -------------------------------------------------------------------------%

Hierarchical task decomposition is the practice of decomposing complex goals into smaller subtasks through a structured hierarchy, and has been central to AI planning since the early work on hierarchical task networks (HTN)~\cite{erol1994htn,erol1996complexity}. In HTN planning, a compound task is decomposed into subtasks via methods, recursively until only primitive actions (operators) remain. The HTN model differs from classical planning in that the search space is guided by domain knowledge encoded as decomposition methods rather than by state-transition heuristics alone. HTN planning theory establishes that hierarchical decomposition is strictly more expressive than classical (STRIPS-style) planning, but becomes undecidable under surprisingly modest conditions when restrictions on non-primitive tasks or task ordering are removed.

Contemporary multi-agent systems have independently rediscovered hierarchical decomposition as a structural necessity for managing long-horizon tasks. ROMA recursively decomposes goals into dependency-aware subtasks, executes them in parallel where possible, and aggregates results to control context growth, standardizing four modular roles---Atomizer, Planner, Executor, and Aggregator---to separate orchestration from model selection and produce transparent, hierarchical execution traces~\cite{roma2025}. ReAcTree builds dynamic agent trees where high-level goals are decomposed into subgoals managed by LLM agent nodes that can reason, act, and further expand the tree, with control-flow nodes coordinating execution strategies~\cite{reactree2025}. StructuredAgent combines online hierarchical planning using dynamic AND/OR trees with structured memory for long-horizon web tasks, producing interpretable plans that facilitate debugging and human intervention~\cite{structuredagent2025}. Deep Agent introduces a Hierarchical Task DAG (HTDAG) that dynamically decomposes high-level goals into interdependent subtasks with a recursive two-stage planner-executor enabling continuous task refinement~\cite{deepagent2025}. MetaGen dynamically adapts roles and collaboration topology at inference time without updating base model weights, effectively managing agent lifecycles through query-conditioned role specifications and constrained execution graphs~\cite{metag2025}.

A critical gap in both classical HTN and contemporary multi-agent decomposition approaches is that decomposition decisions are not recorded as first-class authorization artifacts, and child nodes are not distinguished from parent nodes in terms of their accountability status. A classical HTN planner's subtasks are subroutine calls without independent accountability chains; an LLM agent's spawned sub-agents are independent processes, but the chain of authorization is implicit in the context window rather than explicit in an immutable warrant structure. Dynamic hierarchical task decomposition in open-loop agentic systems has been identified as a vector for \emph{autonomous drift}: the progressive expansion of agent behavior beyond authorized scope because decomposition is unconstrained~\cite{autoGPT2023}. RSP adopts the hierarchical decomposition paradigm but adds the accountability layer absent from both traditions: each node in the decomposition tree is a first-class accountability entity with an immutable parent pointer, a defined termination condition, and a cryptographically verifiable chain back to a human anchor. Decomposition is not merely a planning technique; it is an authorization event that creates a new accountable entity.

% -------------------------------------------------------------------------%
\subsection{The BX3 Framework as Substrate for Immutable Parent Chains}
% -------------------------------------------------------------------------%

The BX3 Framework provides the minimal sufficient three-layer accountability architecture on which the Recursive Spawning Protocol operates~\cite{bx3-framework,bx3framework2026}. BX3 defines three functional layers---the \textit{Purpose Layer}, the \textit{Bounds Engine}, and the \textit{Fact Layer}---each defined by the properties it must maintain rather than by the type of actor occupying it. The Purpose Layer carries intent, judgment, and final human accountability; the Bounds Engine carries constrained execution within a defined operating range; the Fact Layer carries deterministic enforcement and forensic auditability. RSP is the second named pillar of the BX3 Framework, complementing Loop Isolation, Spatial Firewall, Sandbox Gate, and the Bailout Protocol~\cite{bx3-framework,bansal2019}.

The upstream accountability guarantee of the BX3 Framework states that when any node encounters a state it cannot resolve within its provisioned bounds, accountability escalates recursively upward through the system hierarchy, bypassing all machine actors, until it reaches a human anchor~\cite{bx3-framework,bx3framework2026}. This guarantee is not a policy choice; it is a structural property of the three-layer architecture. The Fact Layer creates the immutable parent pointer at spawn time, issues the spawn warrant, and maintains the forensic ledger that records every spawn event. The Fact Layer's determinism is what makes the parent pointer architecturally immutable: once written, the pointer cannot be updated because the Fact Layer does not permit modifications to its authorization records.

The depth bound $D_{\max}$ is a Purpose-Layer-defined, Bounds-Engine-enforced, Fact-Layer-hardened constraint. The Purpose Layer defines the maximum permissible chain depth as part of the spawn authorization policy. The Bounds Engine evaluates each proposed spawn against this limit, returning \texttt{SPAWN\_REFUSED} if the spawn would exceed it. The Fact Layer's pre-commit hook enforces this limit structurally: any spawn operation that would produce $|C| \geq D_{\max}$ is rejected before execution, regardless of the Bounds Engine's return. The constraint is thus enforced three times---by the layer that defines it, the layer that evaluates against it, and the layer that structurally blocks violations. Removing any layer causes the corresponding guarantee to collapse: without the Purpose Layer, there is no authoritative human anchor and no source of spawn authorization policy; without the Bounds Engine, depth limits become unenforceable conventions; without the Fact Layer, immutability of the parent pointer and tamper-evidence of the ledger are promises rather than constraints.

The BX3 pillars collectively establish that BX3 systems fail upward into human consciousness, never downward into autonomous chaos. RSP specifically addresses the structural conditions that enable autonomous drift---the mutable parent pointer, the absent scope refinement constraint, and the missing depth bound enforcement---making the drift triad architecturally impossible through the Fact Layer's write protocol immutability, the Purpose Layer's spawn authorization gate, and the Bounds Engine's depth enforcement. The BX3 Framework is not merely compatible with RSP; it is the minimal sufficient architectural substrate on which RSP's guarantees are structurally realizable.